\chapter{First Nouns and the Article}

Greek nouns change their endings to show their job in the sentence. This is
called \emph{declension}. English shows almost none of this, so the idea
feels new; the pattern itself is small and regular.

\section{The second declension: \PSEGreek{λόγος}}

The word \PSEGreek{λόγος} (word, account) shows the endings of a whole family of
masculine nouns:

\begin{table}[ht]
  \centering
  \begin{tabular}{lll}
    \textbf{Case} & \textbf{Singular} & \textbf{Plural} \\
    Nominative & \PSEGreek{λόγος} & \PSEGreek{λόγοι} \\
    Genitive & \PSEGreek{λόγου} & \PSEGreek{λόγων} \\
    Dative & \PSEGreek{λόγῳ} & \PSEGreek{λόγοις} \\
    Accusative & \PSEGreek{λόγον} & \PSEGreek{λόγους} \\
  \end{tabular}
  \caption{Declension of \PSEGreek{λόγος}.}
\end{table}

\section{The article}

Greek has one word for \emph{the} and (much more rarely) \emph{a/an}. The
article agrees with its noun in gender, number, and case; \PSEGreek{ὁ} is masculine,
\PSEGreek{ἡ} feminine, \PSEGreek{τό} neuter. Learning the article's full pattern is the single
highest-value investment of your first month.

\section{Reading: a first sentence}

\PSEGreek{ὁ} \PSEGreek{ἀγαθὸς} \PSEGreek{λόγος} — the good word. In later chapters the endings will tell
you who does what, so word order can be far freer than in
English.\footnote{Modern Greek, descended from this language, still
declines its articles.}

\PSEIndex{declension}
\PSEIndex{second declension}
